% Converted from the authoritative recovery source, 2026-11-14.
\section{Recoverable-source influences: an actual AA restriction
proof}\label{recoverable-source-influences-an-actual-aa-restriction-proof}

14 November 2026. New deductions from the uniform AA inequality and the
fractional-certificate bound recorded in the supplied exposition.
Constants below are absolute unless their structural dependence is
stated. No claim of priority is made.

This note proves a generalized influence inequality and uses it in the
original restriction/total-influence proof. It does not infer influence
progress from an already established simulation theorem. It also
upgrades the existing forest and bounded-degree high-temperature results
from quantum acceptance functions to \textbf{every polynomial bounded on
the entire cube}.

\subsection{Two imported inputs, and a robustness
lemma}\label{influence-recovery:two-imported-inputs-and-a-robustness-lemma}

For a function bounded in \([0,1]\) on the uniform Boolean cube, the
imported AA bound is

\[
 \max_j \mathbb E(D_jg)^2
 \ge c_0\frac{\operatorname{Var}(g)^{5/2}}{\deg(g)^6}.
 \tag{U}
\]

The supplied exposition, Fact 2.4, gives an absolute \(C\) such that
every globally bounded degree-\(d\) polynomial \(p\) has a fractional
certificate at each \(x\): nonnegative weights \(\lambda_x(i)\) of total
mass at most \(Cd^2\) satisfying

\[
 |p(x)-p(y)|\le\sum_{i:x_i\ne y_i}\lambda_x(i).
 \tag{FC}
\]

Consequently, if a coupling \((X,Y)\) obeys
\(\Pr[Y_i\ne x_i\mid X=x]\le\gamma\) for every \(i\) and every supported
\(x\), then

\[
 \mathbb E|p(X)-p(Y)|\le C d^2\gamma.
 \tag{R}
\]

Indeed, condition (FC) on \(X=x\) and sum its weights. The conditional
hypothesis is essential. No independence of the errors is required.
Since both values lie in \([0,1]\), this also bounds their mean squared
difference. In particular,

\[
 \left|\sqrt{\operatorname{Var}p(X)}-
             \sqrt{\operatorname{Var}p(Y)}\right|
 \le\sqrt{C d^2\gamma}.
 \tag{RV}
\]

The last inequality is the \(L^2\) triangle inequality after applying
the orthogonal projection that subtracts the mean.
